%% kapa: exam=1 round=30 year=2019 period=1 form=A subject=real_estate subject_ko=부동산학원론 seq=11 number=91
%% source: Q-Net 감정평가사 기출문제 (cst00309) artlSeq=5205205 file=extracted/1008/5205205/2019년 제30회 감정평가사 1차시험 문제지/2019년 제30회 감정평가사 1차시험 문제지/2019년도 제30회 감정평가사 1차 1교시 A형.hwp
%% conversion: kordoc 4.24.51 -> tools/gen1.py (auto); review=proofread-sonnet; flags=none
\begin{question}{91}
고정금리대출과 변동금리대출에 관한 설명으로 옳은 것은?
\begin{choices}
  \item 예상치 못한 인플레이션이 발생할 경우 대출기관에게 유리한 유형은 고정금리대출이다.
  \item 일반적으로 대출일 기준 시 이자율은 변동금리대출이 고정금리대출보다 높다.
  \item 시장이자율 하락 시 고정금리대출을 실행한 대출기관은 차입자의 조기상환으로 인한 위험이 커진다.
  \item 변동금리대출은 시장상황에 따라 이자율을 변동시킬 수 있으므로 기준금리 외에 가산금리는 별도로 고려하지 않는다.
  \item 변동금리대출의 경우 시장이자율 상승 시 이자율 조정주기가 짧을수록 대출기관에게 불리하다.
\end{choices}
\end{question}
\answer{3}
